\documentclass[11pt]{article}
\usepackage[utf8]{inputenc}	% Para caracteres en español
\usepackage{amsmath,amsthm,amsfonts,amssymb,amscd}
\usepackage{multirow,booktabs}
\usepackage[table]{xcolor}
\usepackage{fullpage}
\usepackage{lastpage}
\usepackage{enumitem}
\usepackage{fancyhdr}
\usepackage{mathrsfs}
\usepackage{wrapfig}
\usepackage{setspace}
\usepackage{calc}
\usepackage{multicol}
\usepackage{cancel}
\usepackage[retainorgcmds]{IEEEtrantools}
\usepackage[numbers,sort&compress]{natbib}
\usepackage[margin=3cm]{geometry}
\usepackage{amsmath}
\newlength{\tabcont}
\setlength{\parindent}{0.0in}
\setlength{\parskip}{0.05in}
\usepackage{empheq}
\usepackage{framed}
\usepackage[most]{tcolorbox}
\usepackage{xcolor}
\colorlet{shadecolor}{orange!15}
\parindent 0in
\parskip 12pt
\geometry{margin=1in, headsep=0.25in}
\theoremstyle{definition}
\newtheorem{defn}{Definition}
\newtheorem{reg}{Rule}
\newtheorem{exer}{Exercise}
\newtheorem{note}{Note}
\newcommand{\printreferences}{%
  \bibliographystyle{unsrtnat}%
  \bibliography{references}%
}
\begin{document}
\setcounter{section}{0}
\title{NOTE}

\thispagestyle{empty}

\begin{center}
{\LARGE \bf NOTE}\\
{\large Yuqing Rong}\\
Fall 2026
\end{center}

\section{Motivation}
application: fault tolerant quantum computing.

connection between phase of matter and the application on quantum computer ?

how to find the error from measurement ?

An TO-trivial phase transition occurs at a critical value of the error rate;
if the error rate is below the critical value (the accuracy threshold), encoded information can be protected arbitrarily well in the limit of a large code block.
find this threshold ?

A measure error (boundary error) comes from multiple possible bulk errors:  $P(m) = \sum_E P(m|E) P(E)$

$E_{virtual} -> E_{physical} -> E(m)$

\section{Method}

\subsection{isoTNS representation}

We consider the toric code on a square lattice with open boundary conditions,
\[
H=-\sum_s A_s-\sum_p B_p, \qquad
A_s=\prod_{i\in s}\sigma_i^z, \qquad
B_p=\prod_{i\in p}\sigma_i^x.
\]

\begin{figure}[htbp]
  \centering
  \includegraphics[width=0.4\textwidth]{image/lattice_diagram.pdf}
  \caption{Square-lattice geometry of the planar code with open boundaries.}
  \label{fig:planar-lattice}
\end{figure}

Each bond is represented by a doubled Hilbert space, \(\lvert i\rangle\mapsto
\lvert ii\rangle\). The local rank-eight tensor is
\[
T_{\alpha\beta\gamma\delta}^{ijkl}
=
\left\{
\begin{aligned}
  \frac{1}{\sqrt{2}}\,
  \delta_{i\alpha}\delta_{j\beta}\delta_{k\gamma}\delta_{l\delta},
  &\quad \text{if } \alpha+\beta+\gamma+\delta=0 \pmod{2}, \\
  0, &\quad \text{otherwise}.
\end{aligned}
\right.
\]
The physical and virtual indices are ordered as east, north, west, and south.

\begin{figure}[htbp]
  \centering
  \includegraphics[width=0.4\textwidth]{image/lattice_diagram_double_Hilbert_space.pdf}
  \caption{Planar-code lattice in the doubled-Hilbert-space representation. Each virtual bond is doubled as \(\lvert i\rangle\mapsto\lvert ii\rangle\), and \(T_{\alpha\beta\gamma\delta}^{ijkl}\) labels the local rank-eight tensor.}
  \label{fig:doubled-planar-lattice}
\end{figure}

Independent bit-flip errors occur on the internal virtual bonds with
probability \(p\). A distance-\(d\) patch contains \(2d(d-1)\) such bonds.

\subsection{Isometric tensors and sequential circuit}

\subsubsection{Local isometries}

The local tensor can be regarded as an isometry and completed to a unitary
gate. This correspondence provides a circuit representation of the tensor
network.

\begin{figure}[htbp]
  \centering
  \includegraphics[width=0.5\textwidth]{image/isometric_constraint.png}
  \caption{Local tensor \(T\) as an isometry and its completion to a unitary gate \(U\). The input legs \(\gamma,\delta\), together with four physical ancillas initialized in \(\lvert0000\rangle\), are mapped to the output legs \(\alpha,\beta,i,j,k,l\); the lower panels show the reshaped tensor and the isometry condition \(T^\dagger T=I\).}
  \label{fig:local-isometry}
\end{figure}

\subsubsection{Boundary conditions and gate schedule}

The planar patch has two rough and two smooth boundaries and encodes one
logical qubit. Physical qubits are initialized in \(\lvert 0\rangle\). The
virtual boundary convention is \(\lvert +\rangle\) on the west and east
boundaries and \(\lvert 0\rangle\) on the south and north boundaries. At the
east and north boundaries, the terminal carrier is absorbed into the local
gate.

There are four local-gate types: a six-qubit bulk gate, five-qubit east and
north boundary gates, and a four-qubit corner gate. Each gate satisfies the
isometry condition. Gates are applied in diagonal layers, progressing from
south to north and from west to east.

\subsection{Sampling and virtual-bond diagnostics}

The circuit admits a classical message-passing description. The Yao.jl
implementation is a small-system reference: it applies each local gate to a
state vector, performs a Born measurement, records the result, and resets the
measured physical qubits. The production estimator instead works directly
with batches of virtual-bond error frames. For each shot, it samples
independent horizontal and vertical bond-error bits with probability \(p\),
then evaluates the detector record and west-boundary logical parity by XOR
operations. The full detector batch is decoded once with PyMatching, and a
logical failure is counted when the predicted logical parity differs from the
sampled one. No unitary gates or individual physical measurement outcomes are
simulated.



\subsubsection{Virtual-bond noise checks}

For an \(L\times L\) lattice, with \(L=d\), the following quantities diagnose
the virtual-bond noise and its percolation behaviour:
\begin{enumerate}
  \item \textit{Injected-error and mismatch densities:}
  \begin{align*}
  \rho_{\mathrm{injected}}
  &=\frac{\text{number of erroneous virtual bonds}}{2L(L-1)+4L}\approx p, \\
  \rho_{\mathrm{mismatch}}
  &=\frac{\text{number of mismatched internal doubled bonds}}{2L(L-1)}\approx p.
  \end{align*}

  \item \textit{Plaquette frustration:} the probability that a plaquette has
  odd error parity is
  \[
  P_{\mathrm{odd}}
  =\binom{4}{1}p(1-p)^3+\binom{4}{3}p^3(1-p)
  =\frac{1}{2}\left[1-(1-2p)^4\right].
  \]

  \item \textit{Largest occupied cluster:} the normalized cluster size is
  \(S_{\max}/L^2\), where \(S_{\max}\) is the number of sites in the largest
  connected error cluster.

  \item \textit{Horizontal spanning probability:}
  \[
  P_{\mathrm{span}}^{\mathrm{horizontal}}
  =P(\text{an error cluster connects the west and east boundaries}).
  \]
\end{enumerate}


\subsection{Syndrome extraction}

A detector is a binary indicator of whether an error violates an expected
parity constraint. Let \(\mathbf e\) denote the binary virtual-bond error
vector. The detector record is
\[
\mathbf s=\mathsf H\mathbf e \pmod{2},
\]
where \(\mathsf H\) is the parity-check matrix. Equivalently,
\[
s_i=
\left\{
\begin{aligned}
  0, &\quad \text{if an even number of adjacent bonds have errors}, \\
  1, &\quad \text{if an odd number of adjacent bonds have errors}.
\end{aligned}
\right.
\]

For bit-flip noise, we use two detector families:
\begin{enumerate}
  \item an elementary plaquette detector,
  \[
  S_{r,c}=h_{r,c}\oplus h_{r+1,c}\oplus v_{r,c}\oplus v_{r,c+1};
  \]
  \item north and south boundary detectors,
  \[
  n_c=h_{1,c}, \qquad s_c=h_{d,c}.
  \]
\end{enumerate}
The total number of detectors is
\((d-1)^2+2(d-1)=d^2-1\).

\subsection{Decoding and logical-failure estimation}

PyMatching returns both a minimum-weight correction and the logical parity
associated with that correction. Let \(E\) be the physical error, \(C\) the
decoder correction, and \(R=E\oplus C\) the residual error. A logical failure
occurs when
\[
l(R)=l(E)\oplus l(C)=1.
\]

The finite-size scaling ansatz for the logical-failure probability is
\[
P_{\mathrm{fail}}(d,p)\simeq F\!\left((p-p_c)d^{1/\nu}\right).
\]
For each \((d,p)\) point and each batch, we:
\begin{enumerate}
  \item sample all \(2d(d-1)\) virtual-bond errors for multiple shots;
  \item compute the \(d^2-1\) detector values algebraically;
  \item compute the hidden west-boundary logical parity;
  \item decode the batch with PyMatching; and
  \item count disagreements between the actual and predicted logical parities.
\end{enumerate}
For \(N\) shots with \(F\) logical failures,
\[
P_{\mathrm{fail}}=\frac{F}{N},
\qquad
\mathrm{SE}(P_{\mathrm{fail}})
=\sqrt{\frac{P_{\mathrm{fail}}(1-P_{\mathrm{fail}})}{N}}.
\]


\subsection{Bootstrap uncertainty}

To quantify the uncertainty of a crossing estimate, we partition the shots at
each \((d,p)\) point into equal batches. Each bootstrap replicate resamples
these batch failure rates with replacement, recomputes the monotone
logical-failure curves for adjacent code distances, and extracts their unique
crossing. The confidence interval is obtained from the corresponding empirical
quantiles of the valid bootstrap crossings.


\newpage
\section{compile with Clifford gate}
To encode the $|0_L \rangle$ state, we propose two ways:

1. based on $|0_L\rangle\propto (I+\bar Z)\prod_s(I+A_s)|+\rangle^{\otimes N}$

2. based on $|0_L\rangle \propto \prod_p (I+ B_p) |0 \rangle^{\otimes N}$

Here is an $3*3$ example:
\begin{figure}[htbp]
  \centering
  \includegraphics[width=0.4\textwidth]{../document/image/3*3lattice.png}
  \caption{$3*3$ lattice. purple: $A_s$, blue: $B_p$}
  \label{fig:planar-lattice}
\end{figure}

1. $A_s$-based encoding:

(a) A CNOT can establish an even-parity constraint.
eg:
\[
Z_2Z_3Z_5Z_6=+1
\quad\Longleftrightarrow\quad
x_2\oplus x_3\oplus x_5\oplus x_6=0.
\]If qubit \(5\) starts in \(|0\rangle\), the CNOTs
\[
2\to5,\qquad3\to5,\qquad6\to5
\]compute \(x_5=x_2\oplus x_3\oplus x_6\), establishing that parity.

(b) choose qubit 2 to fix the logical-zero parity, 4,5,7,9 as fresh parity targets. These qubits start in \(|0\rangle\);

leaves \(1,3,6,8\) as independent inputs, these qubits start in \(|+\rangle\), so we apply
\[
H_1,\ H_3,\ H_6,\ H_8.
\]

(c)  \(\bar Z=Z_1Z_2Z_3\). Apply
\[
1\to2,\qquad3\to2.
\]Since \(x_2\) started at zero, it becomes \(a\oplus b\). Therefore,
\[
x_1\oplus x_2\oplus x_3
=a\oplus(a\oplus b)\oplus b=0.
\]This establishes \(\bar Z=+1\).

(d) Build the four As constraints block by block.

\begin{figure}[htbp]
  \centering
  \includegraphics[width=0.6\textwidth]{../results/rotated-planar-native-yao-d3-d7-2026-09-12/as-d3-zero-x_ns.pdf}
  \caption{$A_s$-based encoding for $3*3$ lattice.}
  \label{fig:planar-lattice}
\end{figure}

2. $B_p$-based encoding:

choose represensitive qubits for every $B_p$ plaquette, and apply $H$, others are target qubits.

\begin{figure}[htbp]
  \centering
  \includegraphics[width=0.6\textwidth]{../results/rotated-planar-native-yao-d3-d7-2026-09-12/bp-d3-zero-x_ns.pdf}
  \caption{$B_p$-based encoding for $3*3$ lattice.}
  \label{fig:planar-lattice}
\end{figure}

\newpage
\section{error model}

apply X-only or Z-only error after each elementary gate (H, CNOT) with rate $p$

\section{sampling}

Monte Carlo sampling chooses which faults occur; the Pauli-frame simulator determines their effect on the circuit.

For each code distance \(d\) and physical error probability \(p\):

1. Sample faults. After every gate, independently apply an X fault with probability \(p\) to each participating qubit. Each CNOT therefore has two possible fault locations.

2. Propagate errors. Track the X/Z error pattern through the remaining circuit using exact Clifford rules,
   \[
   E\longrightarrow UEU^\dagger.
   \]This accounts for errors spreading and cancelling, without storing the full wavefunction.

3. Decode and judge success. Calculate the final, ideal stabilizer outcomes and give them to the matching decoder. Compare its proposed correction with the actual propagated error. For preparation of \(|0_L\rangle\), a remaining logical X counts as failure.

4. Repeat and estimate. Run \(N\) independent trials and count \(F\) failures:
   \[
   \boxed{\widehat P_L(d,p)=\frac{F}{N}.}
   \]
Error propagation is exact for the Clifford circuit and Pauli-noise model; the estimated failure probability has finite-sampling uncertainty.

\section{results(2D)}

1. The two Z-only panels

The Z-only noise at most cause Logical-Z faults

\[
\bar Z|0_L\rangle=|0_L\rangle.
\]Logical-Z faults can occur without changing the $|0_L\rangle$ state. Maybe need the logical-Z failure curves?

2. The X-only panels

The As/X is fragile logical-parity preparation. Single X fault during logical-state preparation can become a complete logical X error
eg, in Figure 4, only $X_2$ would cause a whole logical-X string.

In the parity part, for \(m\) independent locations, their odd-flip probability is
\[
q_m(p)=\frac{1-(1-2p)^m}{2},
\]which imposes a logical-failure floor. Increasing \(m\) makes this floor larger, approaching \(1/2\), not smaller.

contain more CNOTs, so more fragile to X-only noise.

\begin{figure}[htbp]
  \centering
  \includegraphics[width=\textwidth]{../results/rotated-planar-plus-input-2026-09-22/critical-window-fit/plus_input_logical_errors-local-fit.pdf
  }
  \caption{X error and Z error. As based and Bp based circuit is just for different basis, they are symmetry. Bp/X: X propagate only form short-range error; Bp/Z: Z propagate form long-range error.}
  \label{fig:planar-lattice}
\end{figure}

\begin{figure}[htbp]
  \centering
  \includegraphics[width=\textwidth]{../results/bp-y-h-threshold-2026-09-23/logical_errors.pdf
  }
  \caption{Y error and (X+Z)/$\sqrt(2)$ error. Y: seperate decoder treat X and Z separately. (X+Z)/$\sqrt(2)$: rotate an angle only change prefactor, the logical Z threshold would always be 0 only if Z element is not zero.}
  \label{fig:planar-lattice}
\end{figure}


\newpage
\section{3D toric code}

\subsection{Geometry, checks, and logical operators}

We use a relative cubic slab with periodic $y,z$ directions and rough open
$x$ ends.  There are $L_x$ retained vertex planes $x=1,\ldots,L_x$;
$x$-edges lie at $x=0,\ldots,L_x$, while transverse edges lie at
$x=1,\ldots,L_x$.  Coordinates $y\in\mathbb Z_{L_y}$ and
$z\in\mathbb Z_{L_z}$ are periodic.  The two periodic length-two cells
are retained as distinct cells, and boundary plaquettes are triangular.
We allow $L_x\geq1$ and $L_y,L_z\geq2$.  With $M=L_yL_z$,
\[
 H=-\sum_v A_v-\sum_f B_f,\qquad
 A_v=\prod_{e\ni v}X_e,\quad B_f=\prod_{e\in\partial f}Z_e,
 \qquad n=(3L_x+1)M.
\]
There are $N_A=L_xM$ independent stars and
$N_B=(3L_x+2)M$ plaquettes of rank $(2L_x+1)M-1$.
Thus $k=n-N_A-\operatorname{rank}B=1$.
For example, the two rough-end $xy$ checks are
\[
 B^{xy}_{0,y,z}=Z^x_{0,y,z}Z^x_{0,y+1,z}Z^y_{1,y,z},\qquad
 B^{xy}_{L_x,y,z}=Z^x_{L_x,y,z}Z^x_{L_x,y+1,z}Z^y_{L_x,y,z}.
\]
The $xz$ end checks follow by exchanging $y$ and $z$; interior plaquettes
have four edges.  The $2\times2\times2$ reference therefore has 28 qubits,
8 independent stars, and 32 plaquettes of rank 19.  Its circuit has
8 Hadamards and 40 CNOTs, with a checkpoint after each four-star plane.
A longitudinal $Z$ string joining the rough ends has distance
$d_Z=L_x+1$; a transverse $X$ membrane has distance $d_X=M$.
The prepared state is the normalized projector state
\[
 |0_L\rangle=2^{N_A/2}\prod_v\frac{I+A_v}{2}|0\rangle^{\otimes n}.
\]

\subsection{Sequential circuit, noise, and feedback}

Visit vertices in increasing $(x,y,z)$ order.  At each vertex the fresh
outgoing $x$-edge is Hadamard transformed, then controls CNOTs onto the
incoming $x$-edge and the $+y,-y,+z,-z$ edges, in that order.  The last
outgoing $x$-edges terminate at the right rough end.  For $L_y=2$ or
$L_z=2$, opposite directions are still distinct edge qubits.  We assume
ideal initialization, idle periods, stabilizer measurement, and correction.
Each active output of an H or CNOT independently acquires a Bernoulli-$p$
$X$ error in the X-only experiment, or a Bernoulli-$p$ $Z$ error in the
Z-only experiment.  The inactive Pauli channel has rate zero.

Both protocols replay the same sampled gate faults.  Final-only recovery
decodes the full syndrome once after preparation.  Layer feedback measures
and decodes after every completed $x$ layer, correcting only qubits in the
current prefix slab.  Its frontier has moving triangular plaquettes; future
qubits are unavailable for correction.  At each checkpoint, a Z recovery
from star defects uses matching and an X recovery from plaquette defects uses
belief propagation with ordered-statistics decoding.  Matching is only the
Z-recovery component, not a decoder for both CSS channels.  The small-size
exact option minimizes data-recovery weight; it is not exact circuit-level
maximum-likelihood inference.  Decoder ties, finite BP iterations, and the
ideal measurement assumption limit conclusions from the curves.

The reported residual logical-X and logical-Z parities are Pauli-frame
diagnostics.  Either-logical failure is their union; state failure for the
prepared $|0_L\rangle$ is logical-X failure alone.  In particular a logical
$Z$ changes a Pauli diagnostic but leaves $|0_L\rangle$ unchanged.

\subsection{An indistinguishable Z-fault pair}

Assume $0\leq p<1/2$ and a decoder with access only to measured stabilizer
syndromes and declared model parameters.  For each transverse position
$(y,z)$ choose two independent Z-fault bits:
$a$ on the left boundary $x$-edge after its first CNOT, and $b$ on the
outgoing $x$-edge after the final CNOT of the first-plane gadget.  At
checkpoint $j$, $a$ alone leaves $Z$ on the $x=0$ edge; $b$ alone has
propagated to the $x=1,\ldots,j$ edges.  These two errors give the same
first-plane star syndrome.  Toggling both adds
$D_j=Z_{x=0}Z_{x=1}\cdots Z_{x=j}$, a full prefix longitudinal string
invisible to every checkpoint syndrome.  At the final checkpoint $D_{L_x}$
is logical $Z$.  The $M=L_yL_z$ pairs are disjoint.  Give a decoder a genie
that reveals every other fault and each pair XOR $a\oplus b$, the syndrome
information, but not the constituent bits.  The remaining ambiguity in
$a$ toggles the final logical-Z parity.  If any XOR is one, its two possible
bit assignments have equal probability, so the parity is unbiased.
All XORs vanish with probability
$t^M$, where $t=(1-p)^2+p^2$.  Conditioned on a zero XOR, the two
assignments $(a,b)=(0,0)$ and $(1,1)$ have probabilities $(1-p)^2/t$ and
$p^2/t$.  Their parity bias is $(1-2p)/t$; across all $M$ pairs it is
$((1-2p)/t)^M$.  Even this genie-aided decoder therefore has error
probability
\[
 P_{Z,\mathrm{diagnostic}}\geq
 \frac{1-t^M[(1-2p)/t]^M}{2}
 =\frac{1-(1-2p)^M}{2}.
\]
This is a lower bound for the logical-Z diagnostic under the specified
gate faults and schedule.  It is not a fidelity bound for $|0_L\rangle$:
partial Z strings can stabilize the prepared state, and logical Z itself
acts trivially on that state.

\subsection{Numerical comparison}

The default scan uses cubic $L=3,5,7$, ten physical rates from zero to
$0.15$, and 1000 paired fault records per point.  It reports both distances,
logical-X, logical-Z, either-logical and state failure, Wilson intervals,
and paired protocol differences with uncertainty.  A timed pilot on the
largest size at $p=0.03$ precedes the full run.  The figure overlays the
analytic Z bound; it is a diagnostic reference, not a threshold fit.

A second scan resolves the original Cartesian x-sweep using 10000 paired
fault records at each point, with
\[
 p\in\{10^{-6},3\times10^{-6},10^{-5},3\times10^{-5},
          10^{-4},3\times10^{-4}\},\qquad L\in\{3,5,7,9\}.
\]
Let $p_{1\%}(L)$ denote the linearly interpolated
physical rate at which the logical-Z diagnostic reaches one percent.  With
layer feedback the measured values are
\[
 \bigl(p_{1\%}(3),p_{1\%}(5),p_{1\%}(7),p_{1\%}(9)\bigr)
 = (1.39\times10^{-4},5.03\times10^{-5},
    2.92\times10^{-5},1.51\times10^{-5}).
\]
A four-size power-law summary gives
$p_{1\%}(L)\propto L^{-1.97}$, with a binomial-bootstrap exponent interval
$1.77$--$2.17$.  Thus the apparent small finite-size crossing moves toward
zero rather than approaching a positive threshold, as required by the
decoder-independent bound above.  For example, already at $p=10^{-5}$ the
feedback diagnostic rises from $0.06\%$ at $L=3$ to $0.73\%$ at $L=9$.

\subsection{Direct vertex-projector H+5 experiment}

The primary new diagonal experiment applies the literal cubic-star projector.
This is a paper-motivated construction, not a reproduction of an explicit
three-dimensional gate circuit supplied by the paper. Set
$t=x+y+z$, $u=x$, $v=y$, with open $t=0,\ldots,T$ and periodic
$(u,v)\in\mathbb Z_{L_u}\times\mathbb Z_{L_v}$. An edge $E[d,t,r]$
ends at $V(t,r)$ and begins at $V(t-1,r-s_d)$, where
$s_1=(1,0)$, $s_2=(0,1)$, $s_3=(0,0)$. The encoded slab has
\[
 n=3(T+1)L_uL_v,\qquad k=2,\qquad
 d_Z=2\min(L_u,L_v),\quad d_X=(T+1)\min(L_u,L_v).
\]
Distinct physical edges remain distinct when a transverse period is two.
For the old vertex $V(t-1,r)$, let $e_d=E[d,t-1,r]$ and
$n_d=E[d,t,r+s_d]$. Apply, in this order,
\[
 H(n_1),\quad
 \mathrm{CX}(n_1,n_2),\quad\mathrm{CX}(n_1,n_3),\quad
 \mathrm{CX}(n_1,e_1),\quad\mathrm{CX}(n_1,e_2),\quad
 \mathrm{CX}(n_1,e_3).
\]
There are six gates and eleven participating outputs per cell. The independent
post-output Bernoulli-$p$ X-only or Z-only noise convention is unchanged;
initialization, idles, boundary operations, measurement, and correction are ideal.

Product zero fixes both logical-Z cycles. Ideally project the independent
lower white X triads and correct their outcomes. After the bulk circuit,
project the upper X triads. Their clean outcomes $b$ are uniform on the even
subspace: draw $L_uL_v-1$ bits and fix the last parity. Let $u$ be the flips
induced by the noisy Z frame and let $a=b\oplus u$ be the observed outcome.
Only $a$ enters the physical boundary controller. With its deterministic
correction map $C$, the relative clean/noisy frame is
\[
 z'=z\oplus C(a)\oplus C(b),\qquad a=b\oplus u.
\]
The map uses physical-edge-ID-ordered breadth-first paths to bulk root
$A_{T-1,0,0}$ after deleting the union of both declared logical-X supports.
This exclusion guarantees zero logical overlap of every open correction path;
acyclicity alone does not. The graph stays connected through direction-three
vertical edges and nonwrapping transverse edges. This is a fixed correction
convention and changes noisy logical-Z diagnostics compared with an arbitrary
uncut tree. The clean $b$ and gate faults remain hidden from all decoders.

Every immutable gate-fault record and clean boundary branch is reused by
four arms. Intermediate feedback occurs only after $t=1,\ldots,T-1$:
\begin{enumerate}
\item \texttt{none}: boundary completion and final full-code recovery only.
\item \texttt{one\_step}: in X, measure the periodic six-edge $Q_t$ checks,
pair point defects by unit-length MWPM, and apply the first and last edge
of each canonical path (a length-one path is applied once). Old-target X
faults are unseen but cannot propagate further. Remaining $Q_t$ defects
receive one contraction, not a repeated moving-frontier transport.
In Z, match the retiring/frontier strip $(W,A_0)$ at $t=1$ or
$(A_{t-2},A_{t-1})$ later. Odd parity receives exactly one future dummy;
even parity receives none. Apply an endpoint edge only at a retiring
endpoint, so every retiring syndrome clears while frontier defects may wait.
Symmetric endpoint trimming is not causal in this strip.
\item \texttt{window\_complete}: restrict recovery to $E_{t-1}\cup E_t$.
X uses BP+OSD on $P_1$ initially and $P_{t-1}\cup P_t$ subsequently.
Z uses ordinary open-boundary MWPM on the retiring/frontier strip.
Both clear all completed checks.
\item \texttt{full\_prefix}: recover on all $E_0,\ldots,E_t$ from the
completed plaquettes or lower and bulk stars. Future data are unavailable.
\end{enumerate}
The arms differ in measurements, residual constraints, decoder, and correction
support; feedback-versus-none differences do not isolate correction range.

All arms use common final BP+OSD X recovery and MWPM Z recovery. The fixed
BP baseline is minimum sum, 50 iterations, scale one, parallel schedule,
one thread, and \texttt{OSD\_CS} order two, with uniform prior
$p_{\rm eff}=\min(1-10^{-12},\max(10^{-12},p))$. PyMatching 2.4.0 and
ldpc 2.4.1 are pinned. Matching first minimizes physical length and then
lexicographically fixes endpoint/path signatures, retaining period-two
parallel edges. These decoders omit circuit-induced correlations.
Every complete recovery must satisfy $Hc=s$; a partial route instead obeys
$s'=s\oplus Hr$ with its causal residual constraint. A syndrome-valid
recovery need not restore the intended logical state.

Score both transverse logical-X bits, both logical-Z bits, and all sixteen
Pauli classes. Either-pair logical X is also $|00_L\rangle$ preparation
failure for the X-only experiment. Logical Z is a state-independent Pauli
diagnostic, not $|00_L\rangle$ infidelity. The four-arm statistics retain
sixteen joint masks per observable, Wilson intervals, and all six paired
contrasts. Atomic checkpoints reproduce complete records, including $b$.
A separate timed pilot precedes a 500-shot scout at $L=3,5,7$. The scout
freezes channel-specific production grids and fit windows before 1000-shot
production. Threshold claims require adjacent-size crossing consistency,
an interior monotone likelihood fit, acceptable goodness of fit, and at
least 80\% usable bootstrap fits. Failed guards leave diagnostic candidates.

% H5_MEASURED_RESULTS_BEGIN
The implementation and analysis are being verified; the H+5 scientific scan
has not yet established a threshold. The historical values below belong to
the separate H+6 experiment and are not substituted for H+5 measurements.
% H5_MEASURED_RESULTS_END

\subsection{Historical plaquette-transfer H+6 Z-sector estimate}


The historical seven-gate plaquette-transfer H+6 circuit avoids the earlier
Cartesian propagation mechanism by using
diagonal slices and transferring from old control qubits to new target
triads.  Every single Z gate-output fault then has final support on at most
three qubits in two adjacent slices.  To define a logical observable, use a
periodic honeycomb toric-code seed instead of fixing a unique product
boundary state.  Apply noisy H gates to its $3L_yL_z$ qubits, then ideally
measure its Z-type compatibility checks and apply X corrections.  These
checks commute with the specified independent seed Z faults, so they prepare
the intended logical $|++\rangle$ code space without removing those faults;
no perfect entangled seed is assumed.  The encoded slab has
\[
 n=3(T+1)L_yL_z,\qquad k=2,\qquad
 d_Z=2\min(L_y,L_z),\qquad
 d_X=(T+1)\min(L_y,L_z).
\]
H+5/H+6 threshold values are not a controlled circuit-only comparison:
the seed preparation, upper boundary instrument, and historical
channel-dependent seed noise also differ. The saved H+6 data and interfaces
remain unchanged.

The transfer retains independent rate-$p$ faults on every active H/CNOT
output and ideal measurements, corrections, and idles.  Final unit-weight
matching on the star syndrome returns a full data recovery, whose syndrome
is checked before scoring either transverse logical-Z winding.  For this
Z-only preparation of $|++\rangle$, a residual logical Z changes the state.
The optional feedback in this older scan also measures Z-type compatibility
checks and applies X corrections.  It is therefore inactive under pure Z
noise, so its result agrees exactly with final-only recovery in this sector.
This inactive operation is distinct from the X-star measurements and Z
corrections used by the later three-strength Z experiment.

The completed scan uses $T=L_y=L_z=L=9,13,17,21,25$ and 10,000 shots per
point, for 330,000 production shots.  A quadratic logistic finite-size fit
with a fixed leading $1/L$ correction gives
\[
 p_{c,Z}\simeq 0.0052\quad\text{(approximately $0.52\%$ per active gate output)}.
\]
Fits with different size cuts, probability windows, and correction choices
give estimates roughly $0.49\%$--$0.54\%$; this is a sensitivity range, not
a confidence interval.  The statistical bootstrap, Wilson intervals, raw
counts, model diagnostics, and reproducible commands are recorded in
\texttt{toric3d-z-threshold.md} in \texttt{docs/superpowers} and the associated results.
This is a decoder-dependent Z-sector estimate for the diagonal circuit;
it neither removes the earlier rough-slab bound nor establishes a joint
X/Z threshold.

\subsection{Historical plaquette-transfer H+6 X-sector estimate}

To test the measurement-feedback mechanism itself, use the same encoded
diagonal geometry but prepare logical $|00\rangle$.  Product zero already
fixes the seed Z compatibility checks and both logical-Z cycles.  Ideally
measure the independent seed X stars and apply a syndrome-valid Z correction.
Under the present convention this preparation has no noisy seed gate: only
the thirteen participating outputs of each seven-gate transfer cell acquire
independent X faults of rate $p$.

Before transfer $t$, measure the six-edge Z compatibility checks on active
slice $t-1$.  Their X-error defects are point excitations on a periodic
triangular matching graph.  Unit-weight MWPM pairs the defects.  The primary
one-step rule moves each endpoint one graph edge toward its partner and maps
this virtual old-slice repair through the fixed transfer to a correction on
the fresh targets.  Thus many X errors are removed before they can act as
controls in later transfers.  The classical matching computation is global;
this implementation is not a constant-depth spatially local channel.

After the final slice, measure all physical Z plaquettes and apply a validated
BP+OSD X recovery using minimum-sum BP, 50 iterations, parallel updates,
scaling factor one, one thread, and \texttt{OSD\_CS} order two.  The decoder
uses a uniform physical-$p$ data prior and does not see the sampled fault
record.  State failure is the union of the two residual logical-X membrane
parities.  Unlike the rough-slab logical-Z diagnostic, this event changes the
prepared $|00\rangle$ state.

The production scan uses
\[
 L\in\{5,7,9,11,13\},\qquad
 p\in\{0.0100,0.0125,0.0150,0.0175,0.0200,0.0225,0.0250\},
\]
with $T=L_y=L_z=L$ and 1000 paired fault records per point, for 35000
unique production shots.  Final-only failure increases with size throughout
this window and has no adjacent-size crossing.  One-step feedback reverses
the size dependence near $p=0.015$.  For example, at $p=0.015$ the feedback
failure rates for $L=7,9,11,13$ are respectively
$0.248,0.233,0.211,0.226$, whereas the corresponding final-only rates are
$0.287,0.399,0.464,0.591$.

The all-size quadratic finite-size fit is rejected by its deviance test.  An
all-size cubic alternative gives candidate $p_c=0.01470$ but narrowly misses
the preset goodness guard.  The least aggressive accepted sensitivity uses
the cubic scaling function and $L\geq7$:
\[
 p_{c,X}=0.01482,\qquad \nu=1.056,\qquad
 p_{\mathrm{deviance}}=0.01996.
\]
Of 1000 parametric-bootstrap refits, 999 pass the numerical checks, giving the
conditional 95\% interval
\[
 p_{c,X}\in[0.01428,0.01543].
\]
Accepted size and probability-window variants give $p_c$ from $0.01476$ to
$0.01511$.  The appropriate conclusion is therefore an X-only threshold of
approximately $1.5\%$ per participating gate output, conditional on the
accepted finite-size model and the ideal-feedback assumptions.  A 200-shot
full-path routing check at $L=13,p=0.015$ gives failure $0.200$, compared with
$0.226$ for one-step routing, but its uncertainty is too large for a precise
routing comparison.

This experiment implements the paper's qualitative diagonal-slice mechanism:
measure stabilizers in a two-dimensional slice, pair point defects, and feed
repairs forward before continuing the sweep.  The encoded boundary, transfer
gates, tie breaks, post-gate noise, and final BP+OSD decoder are explicit
choices made here, so the number is not a reproduction of a threshold quoted
by the paper.  It also does not combine this $|00\rangle$ X-sector result with
the preceding $|++\rangle$ Z-sector estimate into a joint Pauli threshold.
Raw counts, Wilson intervals, rejected fits, bootstrap samples, hashes, and
reproduction commands are recorded in
\texttt{toric3d-x-feedback-threshold.md} in \texttt{docs/superpowers} and the
associated results directory.

\subsection{Historical plaquette-transfer H+6 with three feedback strengths}

To hold the diagonal sweep fixed while comparing feedback policies, a further
paired experiment uses the same encoded diagonal transfer and compares three
protocols.  The
\emph{no-feedback} protocol performs no intermediate correction and invokes
only the common final decoder.  For X errors, the \emph{slice-local} protocol
measures the old active slice's Z compatibility checks before the next
transfer.  It computes complete MWPM paths, retains their endpoint steps, and
applies the propagation image of this partial route to the fresh targets after
the noisy transfer.  This is target feed-forward, rather than a direct data
recovery on the old slice.  Its \emph{full-prefix} counterpart measures all
completed prefix Z plaquettes after every nonfinal transfer and applies a
complete BP+OSD X recovery anywhere in the prepared prefix.

For Z errors, the local and full-prefix protocols first measure both colors of
seed X stars and apply a complete MWPM Z recovery on the seed edges; the
no-feedback protocol omits this seed-time feedback.  After each nonfinal
transfer, the local protocol measures a syndrome-complete moving band of three
X-star check sheets and applies a complete MWPM Z recovery on the newest two
data slices.  A literal one-slice Z window would be invalid for the stated
post-gate noise model, because faults on old CNOT controls can place a defect on
the next older check sheet.  The full-prefix Z protocol instead measures all
completed prefix X stars and may apply its complete MWPM Z recovery anywhere
in the prepared prefix.  Unlike the inactive Z-type compatibility operation
in the older Z-only scan, these X-star measurements detect and correct pure Z
faults.

For each channel and Monte Carlo shot, one immutable fault record is replayed
through all three protocols.  Their eight joint Boolean outcomes therefore
give both protocol marginals and paired differences for the full implemented
policies.  In particular, the Z comparison includes the seed-time feedback
difference described above and does not isolate spatial support alone.  Any
complete intermediate or final decoder recovery that fails to reproduce its
measured syndrome aborts the point.  The deliberately partial X one-step
feed-forward need not clear that syndrome, but its underlying complete MWPM
route is validated before truncation.
The X study prepares logical $|00\rangle$ with no noisy seed gates and scores
either residual logical-X membrane.  The Z study prepares logical
$|++\rangle$ with an independent Z fault after every seed H and scores either
residual logical-Z winding.  Thus each reported channel observable is a state
failure, but the two channel thresholds also differ in their logical input,
seed circuit, fault locations, and decoder.

The local protocol has total check participation
$O(TL_yL_z)$, whereas repeatedly decoding the growing prefix costs
$O(T^2L_yL_z)$.  For X noise, the one-step route can leave residual
active-slice syndrome, while full-prefix BP+OSD clears the complete prefix
syndrome.  For Z noise, the three-sheet local protocol already restores zero
complete-prefix star syndrome at every checkpoint.  Full-prefix Z recovery
therefore changes the available paths and logical-class/tie choices, and need
not outperform the local Z protocol.  The X comparison changes the check and
decoder policy as well as correction support, so a threshold increase cannot
be attributed to support range alone.

For the cubic specialization below, $T=L_y=L_z=L$.  The frozen target grid
requested 1000 paired shots for X at
\[
 L=5,7,9,11,13,\qquad
 p=0.025,0.030,0.040,0.045,0.050,0.055,0.060,0.070,0.080,
\]
and 2000 paired shots for Z at
\[
 L=9,13,17,21,25,\qquad
 p=0.025,0.030,0.0325,0.035,0.0375,0.040,0.045,0.050,0.060,0.070.
\]
Measured runtime sets the final size scope to four complete sizes per channel.
X uses $L=5,7,9,11$ at all nine X probabilities, totaling 36 size-rate points
and 36000 paired fault records.  Z uses $L=9,13,17,21$ at the first seven
probabilities through $p=0.045$, totaling 28 points and 56000 paired fault
records.  This preserves a guard point above the transition and both
endpoint-trimmed window checks; the higher frozen Z rates chiefly probe
saturation and were stopped before fitting on the measured runtime.  The
partial $L=13$ X and $L=25$ Z checkpoints are retained only as resumable
provenance and are excluded from every fit and aggregate.

A deliberate interruption revealed a now-fixed bookkeeping bug: a user
interrupt was temporarily labeled as an invalid recovery, and scalar X replay
could record several outcomes before its batch completion counter advanced.
The original interruption artifacts are archived separately.  The live
checkpoints were repaired from atomic counts or deterministic replay, and the
partial X $L=13$, Z $L=25$, and original-shard Z $L=21,p=0.045$ records are not
analysis inputs.  These interruption records are not decoder failures.

Thresholds are exposed only when every adjacent size pair crosses in the fit
window, the monotone likelihood fit passes numerical and deviance checks, and
at least 80\% of 1000 conditional bootstrap fits are usable.  Conditional
Monte Carlo intervals and the spread among accepted size, window, polynomial,
and fixed-$L^{-1}$-correction sensitivities are reported separately.

The runtime-scoped production study contains 36000 X-channel and 56000
Z-channel paired fault records.  All complete recoveries are syndrome-valid,
and an independent reconstruction of the checkpoints, joint tables,
marginals, paired differences, and final analysis inputs finds zero invalid
recoveries.  Recorded decoder construction, sampling, and three-protocol
replay total 10.734 compute hours.

No primary finite-size fit passes every registered threshold guard.  For X
noise, full-prefix feedback gives adjacent-size crossings at 5.159\%, 5.586\%,
and 5.684\%.  Its quadratic diagnostic candidate is
$p_{\mathrm c}=5.297\%$ with $\nu=1.391$, but it is withheld because the
deviance $54.77/31$ has goodness-of-fit $p=0.00530<0.01$.  Accepted cubic and
endpoint-window sensitivity fits lie between 5.293\% and 5.323\%, while a
fixed-$L^{-1}$ correction gives 6.154\%; these remain model diagnostics because
the primary fit fails.  The present no-feedback and slice-local X fits are
inconclusive in this probability window.

For Z noise, slice-local and full-prefix feedback give diagnostic candidates
of 3.934\% and 3.923\%, respectively.  The slice-local data lack the
$L=13,17$ crossing, the full-prefix data lack the $L=17,21$ crossing, and both
quadratic fits are nonmonotone with goodness-of-fit below $1.5\times10^{-5}$.
They therefore show a transition-like onset near 3.9\% but do not establish an
asymptotic threshold.  At $L=21,p=0.035$, the state-failure rates are 0.7545,
0.2445, and 0.2415 for no feedback, slice-local feedback, and full-prefix
feedback.  The paired full-prefix-minus-none difference is $-0.513$ with a
conservative 95\% interval $[-0.549,-0.475]$, whereas full-prefix minus
slice-local is $-0.003$ with interval $[-0.0146,0.0086]$.  Active Z feedback
is therefore effective at fixed size and noise rate, while these data do not
resolve a threshold difference between the local and prefix policies.

These withheld results do not supersede the separate targeted estimates above:
the earlier low-probability Z final-only study gives 0.517\%, and the earlier
one-step X slice-local study gives 1.482\%.  The probability windows in the
present three-policy comparison do not refit either estimate.  Missing
finite-size crossings likewise do not imply a zero threshold.  Complete raw
counts, Wilson intervals, sensitivity fits, paired statistics, hashes, and
reproduction commands are recorded in
\texttt{toric3d-diagonal-feedback-strength.md} in
\texttt{docs/superpowers} and its associated results directory.

\newpage
\section{questions}

\subsection{present 2D toric-code}

Can a different unitary extension reduce the propagation of Z errors? Can additional resources overcome the 2D obstruction?

introduce majority-like correction in 2D:

(a) redesigned unitary schedule (changes plaquette order, pivot selection, CNOT directions, CNOT order.)

(b) repeated stabilizer measurement rounds + classical decoding and feedback (measurement feedback)

(c) redundant encoding of virtual information and correct it coherently


\subsection{move to 3D toric-code}

Use measurement feedback for 3D toric code
\[
\text{prepare one 2D layer}
\rightarrow
\text{measure the layer syndrome}
\rightarrow
\text{decode}
\rightarrow
\text{feedback}
\rightarrow
\text{prepare the next layer}.
\]
compare:

(a) naive 3D sequential circuit

(b) feedback-stabilized 3D sequential circuit

(c) logical failure rate versus noise and system size

\end{document}
